% !PRINTED 383
\weberpage{383}

%% [running head, not body text -- commented out 2026-08-07] Page furniture (running
%% head, folio, signature) is excluded from this edition by standing policy.
%% \noindent\S.\,124. \hfill \textsc{Reducirte Zahlen.} \hfill 383

%<<CONT
(2) $\xi_1^2 = \tfrac{1}{4}$, $\xi_1 = \mp\tfrac{1}{2}$, which is contained both under case a) and under
case b);

and with this, then, Theorem 2. is proved. To this we now add, as a
third theorem:

\noindent 3.\quad \gesperrt{For a given negative discriminant $D$ there is only a finite number of reduced irrational numbers.}

From the representation of the irrational numbers \S.\,122, (7):
\[
\omega = \dfrac{\sqrt{D}+b}{2c}, \qquad \xi = \dfrac{b}{2c}, \qquad \eta = \dfrac{\sqrt{-D}}{2c}
\]
and from the inequality (3) $\eta^2 \geqq \tfrac{3}{4}$ it follows that
\[
\dfrac{-D}{4c^2} \geqq \tfrac{3}{4},
\]
or
\[
c \leqq \sqrt{\dfrac{-D}{3}};
\]
thus, when $D$ is given, there is only a finite number of values
of the positive integer $c$. Further, from $-\tfrac{1}{2} \leqq \xi \leqq \tfrac{1}{2}$ it follows that:
\[
-c \leqq b \leqq c.
\]
Thus for each value of $c$ there is only a finite number of values
for $b$, q.\,e.\,d.

If one unites all numbers equivalent to one another into a
number class, then by 1. we have accordingly proved:

\noindent 4.\quad \gesperrt{The number of classes of imaginary quadratic irrational numbers of a given discriminant is finite.}

\begin{center}
\S.\,124.

\gesperrt{Reduced numbers with positive discriminant.}
\end{center}

We now have to carry out a similar investigation for a
positive $d$, that is, for \gesperrt{real quadratic irrational numbers}. Such a
number $\omega = x + y\sqrt{d}$ shall be called \gesperrt{reduced} if it satisfies the following
conditions:
%CONT>>
